% Main Overleaf-ready LaTeX file containing Chapter 4: Methodology
\documentclass[11pt]{report}
\usepackage[utf8]{inputenc}
\usepackage{graphicx}
\usepackage{booktabs}
\usepackage{hyperref}
\usepackage{amsmath}
\usepackage{longtable}
\usepackage{setspace}
\usepackage{caption}
\usepackage{float}
\usepackage{url}
\usepackage{xurl}
\usepackage{geometry}
\geometry{margin=1in}
\graphicspath{{docs/figures/}}
\hypersetup{colorlinks=true,linkcolor=blue,citecolor=blue,urlcolor=blue}
\usepackage{parskip}
\usepackage{titlesec}
\usepackage{enumitem}
\usepackage{microtype}

\onehalfspacing
\setlength{\parindent}{0pt}
\setlength{\parskip}{6pt}

\titlespacing*{\section}{0pt}{2em}{1em}
\titlespacing*{\subsection}{0pt}{1.5em}{0.8em}
\titlespacing*{\subsubsection}{0pt}{1.2em}{0.6em}

\setlist[itemize]{itemsep=6pt, topsep=4pt}
\setlist[enumerate]{itemsep=6pt, topsep=4pt}

\onehalfspacing
\setlength{\parskip}{0.6em}
\setlength{\parindent}{0pt}

\begin{document}
\onehalfspacing

\title{Methodology: Deepfake Detection System}
\author{Auto-generated from repository}
\date{\today}
\maketitle
\tableofcontents
\clearpage

\chapter{Methodology}

\section{Methodology Overview}
This chapter details the methodology employed to construct, train and evaluate
the deepfake detection system implemented in this repository. The system
integrates multiple complementary signal domains: spatial RGB appearance,
handcrafted high-frequency residuals (\textbf{Steganalysis Rich Model (SRM)}), and frequency-domain magnitude
features (\textbf{Fast Fourier Transform (FFT)}). These modalities are fused
with modern deep feature extractors and lightweight attention/pooling modules.
The pipeline comprises: dataset construction; preprocessing and augmentation;
feature extraction using pretrained backbones; optional preprocessing modules
(\texttt{SRMLayer} and \texttt{FFTLayer}); attention and pooling; optimisation and
training; and evaluation. Implementation details are referenced directly to
repository modules.

The overarching design objective is to produce a classifier that reliably
discriminates among three semantic classes: Real, AI Generated, and AI Edited.
To this end, the methodology emphasises careful dataset curation to avoid
leakage and ensure representativeness, multi-domain preprocessing that
amplifies forensic cues, transfer learning from ImageNet-pretrained backbones
to leverage robust low-level features, and attention and pooling mechanisms
that concentrate model capacity on informative spatial and channel components.
The following sections expand each part of the pipeline while preserving all
factual implementation details present in the codebase.

\section{Dataset Description}
\subsection{Motivation for a Curated Dataset}
Supervised forensic models are sensitive to dataset composition. Publicly
available collections may be biased towards a narrow set of generative models
or capture conditions; therefore the project implements a purpose-built,
multi-source dataset assembled by the pipeline.
This curation enables deterministic pre-processing, enforced class-balancing,
and provenance metadata that supports reproducible experiments and rigorous
error analysis.

\subsection{Three-class Taxonomy}
The dataset uses an explicit three-class taxonomy:
\begin{itemize}
  \item \textbf{Real (Class 0)} — original, unmanipulated images.
  \item \textbf{AI Generated (Class 1)} — fully synthetic images produced by
  generative models (GANs, diffusion models, etc.).
  \item \textbf{AI Edited (Class 2)} — real images that have been edited (for
  example via inpainting, face-swap, or localized retouching).
\end{itemize}

The separation of generated and edited classes enables the classifier to learn
distinct signatures: global spectral fingerprints often indicate fully
generated imagery, while localized residual anomalies better characterise
editing operations.

\subsection{Dataset Size, Splits and Export}
Recorded counts are: real = 26,000; ai_generated = 26,000; ai_edited = 25,865; total = 77,865. The pipeline writes exported splits to export directories. While nominal YAML-configured ratios exist (70/15/15), the implemented cluster-aware splitter produces an empirical export approximating 40/30/30 in some runs; the split implementation is in the splitter module and is orchestrated by the dataset pipeline.

\subsection{Class Balancing and Quota Sampling}
The production pipeline enforces per-source quotas and overall class balance to
prevent dominance of any single class or source. The sampling logic is implemented in the sampler module and deterministic
behaviour is ensured via \texttt{set_global_seed} in the dataset pipeline.
Balanced datasets reduce training bias and ensure that per-class metrics are
meaningful for operational evaluation.

\subsection{Source Diversity and Manipulation Types}
The dataset aggregates diverse capture conditions and manipulation methods to
improve generalisation. Diversity includes multiple camera models, compression
artifacts, and generative model variants; manipulation types include full
generation, inpainting, and face-swapping. The pipeline preserves provenance
metadata (for example an export index file) to support stratified analyses
by source or manipulation family.

\subsection{Deduplication, Validation and Dataset Quality}
Deduplication is performed with perceptual hashing and validation stages check
resolution, blur and metadata (see the deduplicator and validator modules).
Cluster-aware splitting ensures that related items (e.g. frames extracted from the same video) remain
in the same split to prevent leakage. These steps are critical because low-
quality or duplicated training examples encourage memorisation of source
artifacts rather than learning transferable forensic cues.

\section{Data Preprocessing}
\subsection{Overview}
Preprocessing responsibilities are split across dataset-level filtering and
model-input preparation. Model-input preprocessing is implemented primarily in
the preprocessing module and uses Albumentations for training-time
augmentations and Torchvision transforms for deterministic inference-time
processing. Optional preprocessing modules (\texttt{SRMLayer}, \texttt{FFTLayer})
are implemented in dedicated SRM and FFT modules and are applied via the
\texttt{PreprocessNet} wrapper in the training module.

\subsection{Image Loading and Channel Ordering}
At data ingest, images are read with OpenCV (BGR format) within the dataloader
module. The loader performs an explicit BGR→RGB conversion prior to
augmentation to align with pretrained backbones' expected channel ordering and
avoid inadvertent color-space anomalies.

\subsection{Resizing and Normalization}
All inputs are resized to 224×224 during augmentation. Pixel normalisation uses
ImageNet mean and standard deviation (mean=[0.485,0.456,0.406],
std=[0.229,0.224,0.225]) to maintain consistency with pretrained feature
initialisations. Conversion to PyTorch channel ordering (C,H,W) and floating
tensor is performed by \texttt{ToTensorV2} as the final augmentation step.

\subsection{Augmentation Modes and Techniques}
Training augmentations are exposed as three named transforms via
\texttt{get_train_transform(mode)}:
\begin{itemize}
  \item \texttt{light}: Resize(224,224), HorizontalFlip(p=0.5), Normalize(),
  \texttt{ToTensorV2}.
  \item \texttt{standard} (default): Resize(224,224), HorizontalFlip(p=0.5),
  Rotate(limit=15,p=0.5), RandomBrightnessContrast(p=0.5),
  ImageCompression(quality_range=(30,100),p=0.4), Normalize(), \texttt{ToTensorV2}.
  \item \texttt{strong}: as \texttt{standard} with intensified ImageCompression
  (20–95), additive noise (GaussNoise or ISONoise) and light blur (Gaussian or
  Median) to simulate harsh capture conditions.
\end{itemize}

Rationale: Compression simulation addresses lossy storage/transmission; noise
injection models sensor or post-processing noise; geometric transforms improve
robustness to framing variations; blur reduces sensitivity to fine-grained
texture so models learn coarser but more generalisable cues.

\subsection{Training vs Inference Preprocessing}
Training uses stochastic Albumentations to promote invariance, while
inference applies deterministic Torchvision transforms (Resize, CenterCrop,
ToTensor, Normalize) to ensure repeatable inputs and stable evaluation.

\section{Feature Extraction}
\subsection{Candidate Backbones and Feature Dimensions}
The codebase supports multiple backbone families (constructed in the training module): ResNet-18/ResNet-50 (final features
512/2048), ConvNeXt-Tiny/Small (final features 768), EfficientNet-B3 and
ViT-B/16 (supported with special handling for token pooling). Using
ImageNet-pretrained weights provides robust low-level filters and accelerates
convergence when fine-tuning for forensic tasks.

\subsection{Convolutional Inductive Biases for Forensics}
Convolutional architectures are well-suited for forensic analysis due to their
local receptive fields and hierarchical composition. Early layers capture
edges and textures while deeper layers encode increasingly semantic
representations; this hierarchy supports detection of both local manipulations
(inpainting seams) and more global inconsistencies (colour mismatches or
artificial textures). Transformer-based backbones are supported to compare
their global attention properties, but their tokenisation implies different
trade-offs.

\subsection{Pooling and Classification Head}
After spatial aggregation (GeM or average pooling) the pooled vector is
subjected to dropout (default 0.4) and then passed to a linear classification
head mapping to three logits. The head construction is performed by
\texttt{make_head(in_features)} in the training module and
is tuned to each backbone's final feature dimensionality.

\section{Preprocessing Feature Modules}
\subsection{\textbf{Steganalysis Rich Model (SRM)}}
\textbf{Implementation:} \texttt{SRMLayer} (SRM implementation)
computes three fixed high-pass residual channels using handcrafted kernels
(Laplacian 5×5, horizontal 5×5, vertical 5×5) applied to a grayscale image;
residuals are clamped (default clamp range 2.0) and concatenated with RGB to
form a 6-channel input. The helper \texttt{adapt_conv1_for_srm} modifies the
backbone's first convolution to accept 6 channels by initialising additional
weights to 0.1× the pretrained RGB weights.

Justification: SRM residuals amplify small-scale discontinuities and blending
seams that are often obscured by colour and texture in the RGB domain. By
presenting these residuals explicitly to the network, early convolutions can
detect structural anomalies that are indicative of manipulation.

\subsection{\textbf{Fast Fourier Transform (FFT)}}
\textbf{Implementation:} \texttt{FFTLayer} (FFT implementation)
computes the log-magnitude of the centred 2D Fourier transform of the
grayscale image, normalises it to [0,1], and returns a B×1×H×W tensor. When
enabled, this channel is concatenated after RGB and SRM channels (rgb →
residuals → fft) and the backbone's first convolution is adapted accordingly.

Justification: Frequency-domain analysis reveals global periodicities and
power-spectrum anomalies often characteristic of generative models or
resampling introduced by editing tools. FFT magnitude provides complementary
evidence to SRM and RGB features and can distinguish manipulation classes
that are ambiguous in the spatial domain.

\section{Attention and Pooling Mechanisms}
\subsection{\textbf{CBAM Attention}: Channel and Spatial Attention}
CBAM is implemented in the attention_heads module and applies
channel attention (via pooled statistics and an MLP) followed by spatial
attention (via concatenated avg/max maps and a convolution). Applied to late
feature maps, CBAM reweights channels and spatial locations to emphasise
forensically salient information.

\subsection{\textbf{GeM Pooling}}
\textbf{Implementation:} \texttt{GeM} (attention_heads implementation).
Given activation map $X \in \mathbb{R}^{C\times H\times W}$ GeM computes:
\[
\mathrm{GeM}(X)_c =
\left(
\frac{1}{HW}
\sum_{i=1}^{H}
\sum_{j=1}^{W}
X_{c,i,j}^{p}
\right)^{1/p}
\]
The parameter $p$ (default 3.0) is registered as a tensor and may be frozen
or learned. GeM interpolates between average ($p=1$) and max pooling
($p\to\infty$), giving stronger emphasis to local high activations which are
useful when small regions contain manipulation evidence.

\subsection{Role of Attention and GeM in Forensics}
Attention reduces the influence of irrelevant background structure and lets the
model focus on high-signal regions. GeM ensures that strong local cues are
preserved during aggregation rather than being diluted by large spatial
averages; this property is particularly beneficial when manipulations occupy a
small fraction of the image.

\section{Model Architecture}
\subsection{Input and \texttt{PreprocessNet} Wrapper}
The input pipeline optionally applies \texttt{SRMLayer} and/or \texttt{FFTLayer}
to an RGB tensor and concatenates channels in the order rgb → residuals → fft.
The concatenated tensor is provided to an adapted backbone via the
\texttt{PreprocessNet} wrapper implemented in the training module.

\subsection{Backbone, Attention and Pooling}
The adapted backbone (ResNet/ConvNeXt/EfficientNet/ViT) computes hierarchical
features. \textbf{CBAM Attention} blocks may be inserted after late stages
(for ResNet, after \texttt{layer4}; for ConvNeXt, after \texttt{features}) and
\textbf{GeM Pooling} may replace the global average pooling layer. The
resulting pooled vector is passed through a dropout layer and a linear
classification head (softmax for probabilities).

\subsection{Step-by-step Data Flow}
\begin{enumerate}[label=\textbf{Step \arabic*:}]
  \item Image file loaded (cv2) and converted BGR→RGB.
  \item Training augmentations applied (Resize → geometric → photometric →
  Normalize → \texttt{ToTensorV2}) or deterministic inference transforms used.
  \item \texttt{SRMLayer} / \texttt{FFTLayer} modules compute additional
  channels if enabled and outputs are concatenated with RGB.
  \item The concatenated tensor is forwarded through the adapted backbone, with
  optional \textbf{CBAM Attention} applied to late feature maps.
  \item \textbf{GeM Pooling} or average pooling aggregates spatial maps to a
  vector, followed by dropout and the final linear head producing three logits.
\end{enumerate}

\section{Training Procedure}
\subsection{Losses and Class-weighting}
The training code implements CrossEntropy, class-weighted CrossEntropy,
\textbf{Focal Loss} (\texttt{FocalLoss}) and a factory \texttt{build_criterion}
for WeightedFocalLoss. Default class weights are \texttt{[1.5, 1.0, 1.5]}. The
default focal gamma in code is 3.0; however model cards and many experiments
use $\gamma=2.0$ in practice. Label smoothing (commonly 0.1) is supported to
reduce overconfidence.

Rationale: Focal loss focuses optimisation on hard examples, which is useful
when manipulated examples are subtle. Class weights compensate for residual
sampling imbalance and emphasize operationally critical classes.

\subsection{Optimisers, Schedules and Hyperparameters}
Optimisers: Adam (default) and SGD (option). Default weight decay: 1e-4. The
learning rate default is 1e-4. Schedulers: \textbf{Cosine Annealing}
restarts (CosineAnnealingWarmRestarts with \texttt{T_0=20}, \texttt{T_mult=2}) for \texttt{lr_schedule='cosine'}, or ReduceLROnPlateau for \texttt{lr_schedule='plateau'}.
Batch size default is 64 and default epochs is 50 (baseline scripts use
smaller defaults for quick iterations). Early stopping and checkpointing save
the best model (for example to the best model checkpoint) and per-epoch metrics
are written to metrics CSV and training summary JSON files.

Rationale: Cosine restarts provide periodic learning rate resets that help
escape local minima during long runs; Adam provides adaptive gradient steps that
are robust to noisy gradients introduced by augmentations.

\subsection{Mixed Precision and Performance}
Training uses AMP (\texttt{torch.amp.GradScaler}, \texttt{autocast}) and
enables \texttt{torch.compile} when available. These pragmatics increase
throughput and reduce memory footprint, enabling larger batch sizes and
faster experimentation.

\section{Evaluation Metrics}
\subsection{Metrics Implemented}
\begin{itemize}
  \item \textbf{Accuracy}: overall fraction correct; logged as
  \texttt{train_acc} and \texttt{val_acc}.
  \item \textbf{Precision / Recall}: computed per-class; precision measures the
  correctness of positive predictions while recall measures completeness of
  detection for the true positives.
  \item \textbf{F1 score}: harmonic mean of precision and recall; code computes
  macro and per-class F1 using \texttt{sklearn.metrics.f1_score}.
  \item \textbf{Confusion matrix}: reconstructed from saved logits and labels to
  provide detailed misclassification patterns.
\end{itemize}

\subsection{ROC-AUC and Offline Analysis}
ROC-AUC is not computed by default within the training scripts but can be
derived offline from saved logits and labels. ROC-AUC is threshold-independent
and especially useful when selecting operational thresholds that balance false
positives and false-negatives.

\subsection{Complementary Metric Usage}
Accuracy gives a single-number summary but can conceal class-specific failures;
precision, recall and F1 provide per-class diagnostic insight, while the
confusion matrix reveals asymmetric confusions that guide dataset curation and
loss reweighting.

\section{Methodology Summary}
This chapter has presented a principled methodology integrating dataset
engineering, robust preprocessing, multi-domain feature extraction, attention
and pooling mechanisms, and carefully chosen optimisation strategies. The
approach is implemented in the repository with explicit modules to ensure
reproducibility. Each design decision is motivated by forensic considerations:
\textbf{Steganalysis Rich Model (SRM)} emphasises high-frequency residuals for
local edit detection; \textbf{Fast Fourier Transform (FFT)} reveals global
spectral fingerprints; \textbf{GeM Pooling} and \textbf{CBAM Attention}
concentrate model capacity on salient, high-signal regions; and augmentation
plus class-aware sampling promote robust generalisation. The subsequent
chapters evaluate this methodology empirically and discuss limitations and
avenues for future work.

\section{Diagrams}
The following subsections describe recommended diagrams. Each subsection lists
\textbf{Components} and a short \textbf{Purpose} note.

\subsection{Overall System Pipeline}
\textbf{Components}
\begin{itemize}
  \item raw sources
  \item indexing
  \item validation
  \item deduplication
  \item sampling / balancing
  \item cluster-aware splitter
  \item exported train / val / test
  \item training (\texttt{PreprocessNet} → Backbone → Head)
  \item evaluation
  \item model cards / metrics
\end{itemize}
\textbf{Purpose}
End-to-end reproducible workflow from raw data to published experiment artifacts.
 
\subsection{Dataset Preprocessing Pipeline}
\textbf{Components}
\begin{itemize}
  \item indexing
  \item validation checks (resolution / blur)
  \item deduplication (pHash)
  \item quota sampling
  \item cluster-aware split
  \item export
\end{itemize}
\textbf{Purpose}
Clarifies steps that ensure dataset quality and avoid leakage; show conditional branches (invalid → rejected) and deterministic seeding.

\subsection{CNN Backbone Architecture}
\textbf{Components}
\begin{itemize}
  \item input
  \item stem convs
  \item residual / ConvNeXt blocks
  \item stage outputs
  \item pooling (Avg / GeM)
  \item dropout
  \item FC head
\end{itemize}
\textbf{Purpose}
Hierarchical feature extraction for forensic signals (feature sizes: ResNet-18: 512; ConvNeXt: 768; ResNet-50: 2048).

\subsection{SRM Filter Integration}
\textbf{Components}
\begin{itemize}
  \item RGB → grayscale
  \item three SRM kernels (Laplacian, horizontal, vertical)
  \item residual maps
  \item concatenation with RGB
  \item adapted conv1 (0.1× init for extra channels)
\end{itemize}
\textbf{Purpose}
Shows how residual channels are fused at model input and highlights how SRM buffers are stored in checkpoints.

\subsection{FFT Feature Extraction}
\textbf{Components}
\begin{itemize}
  \item grayscale → FFT → fftshift → log-magnitude
  \item normalise → FFT channel concatenated with RGB / SRM
\end{itemize}
\textbf{Purpose}
Documents spectral fingerprint extraction and the transformation from spatial to frequency domain and back to a spatially-aligned magnitude image.

\subsection{Attention Module (CBAM)}
\textbf{Components}
\begin{itemize}
  \item input feature map
  \item channel attention (avg/max pools → MLP → sigmoid)
  \item multiply
  \item spatial attention (concat avg/max → conv → sigmoid)
  \item multiply
\end{itemize}
\textbf{Purpose}
Describes sequential application (channel → spatial) and how attention reweights channels and locations.

\subsection{Feature Extraction Pipeline}
\textbf{Components}
\begin{itemize}
  \item RGB → (optional SRM / FFT)
  \item adapted backbone
  \item attention
  \item pooling (GeM)
  \item feature vector → classifier head
\end{itemize}
\textbf{Purpose}
Shows multi-domain fusion and final prediction flow.

\subsection{Full Model Architecture}
\textbf{Components}
\begin{itemize}
  \item \texttt{PreprocessNet} wrapper
  \item backbone
  \item optional \textbf{CBAM Attention}
  \item \textbf{GeM Pooling}
  \item dropout
  \item fully connected classifier
\end{itemize}
\textbf{Purpose}
Full neural network inference pipeline with tensor dimension and feature flow annotations.

\subsection{Training Workflow}
\textbf{Components}
\begin{itemize}
  \item DataLoader
  \item augmentations
  \item forward pass
  \item loss computation
  \item backward pass
  \item optimizer step
  \item scheduler update
  \item checkpoint saving
\end{itemize}
\textbf{Purpose}
Describes the training loop execution.

\subsection{Evaluation Pipeline}
\textbf{Components}
\begin{itemize}
  \item test dataset
  \item inference
  \item logits collection
  \item metrics computation (accuracy, precision, recall, F1)
  \item confusion matrix generation
\end{itemize}
\textbf{Purpose}
Evaluation and result reporting.

\end{document}
\usepackage{graphicx}
\usepackage{booktabs}
\usepackage{hyperref}
\usepackage{amsmath}
\usepackage{longtable}
\usepackage{setspace}
\usepackage{caption}
\usepackage{float}
\usepackage{url}
\usepackage{xurl}
\usepackage{geometry}
\geometry{margin=1in}
\graphicspath{{docs/figures/}}
\hypersetup{colorlinks=true,linkcolor=blue,citecolor=blue,urlcolor=blue}
\usepackage{parskip}
\usepackage{titlesec}
\usepackage{enumitem}
\usepackage{microtype}

\onehalfspacing
\setlength{\parindent}{0pt}
\setlength{\parskip}{6pt}

\titlespacing*{\section}{0pt}{2em}{1em}
\titlespacing*{\subsection}{0pt}{1.5em}{0.8em}
\titlespacing*{\subsubsection}{0pt}{1.2em}{0.6em}

\setlist[itemize]{itemsep=6pt, topsep=4pt}
\setlist[enumerate]{itemsep=6pt, topsep=4pt}

\onehalfspacing
\setlength{\parskip}{0.6em}
\setlength{\parindent}{0pt}

\begin{document}
\onehalfspacing

\title{Methodology: Deepfake Detection System}
\author{Auto-generated from repository}
\date{\today}
\maketitle
\tableofcontents
\clearpage

\chapter{Methodology}

\section{Methodology Overview}
This chapter details the methodology employed to construct, train and evaluate
the deepfake detection system implemented in this repository. The system
integrates multiple complementary signal domains: spatial RGB appearance,
handcrafted high-frequency residuals (\textbf{Steganalysis Rich Model (SRM)}), and frequency-domain magnitude
features (\textbf{Fast Fourier Transform (FFT)}). These modalities are fused
with modern deep feature extractors and lightweight attention/pooling modules.
The pipeline comprises: dataset construction; preprocessing and augmentation;
feature extraction using pretrained backbones; optional preprocessing modules
(\texttt{SRMLayer} and \texttt{FFTLayer}); attention and pooling; optimisation and
training; and evaluation. Implementation details are referenced directly to
repository modules.

The overarching design objective is to produce a classifier that reliably
discriminates among three semantic classes: Real, AI Generated, and AI Edited.
To this end, the methodology emphasises careful dataset curation to avoid
leakage and ensure representativeness, multi-domain preprocessing that
amplifies forensic cues, transfer learning from ImageNet-pretrained backbones
to leverage robust low-level features, and attention and pooling mechanisms
that concentrate model capacity on informative spatial and channel components.
The following sections expand each part of the pipeline while preserving all
factual implementation details present in the codebase.

\section{Dataset Description}
\subsection{Motivation for a Curated Dataset}
Supervised forensic models are sensitive to dataset composition. Publicly
available collections may be biased towards a narrow set of generative models
or capture conditions; therefore the project implements a purpose-built,
multi-source dataset assembled by the pipeline.
This curation enables deterministic pre-processing, enforced class-balancing,
and provenance metadata that supports reproducible experiments and rigorous
error analysis.

\subsection{Three-class Taxonomy}
The dataset uses an explicit three-class taxonomy:
\begin{itemize}
  \item \textbf{Real (Class 0)} — original, unmanipulated images.
  \item \textbf{AI Generated (Class 1)} — fully synthetic images produced by
  generative models (GANs, diffusion models, etc.).
  \item \textbf{AI Edited (Class 2)} — real images that have been edited (for
  example via inpainting, face-swap, or localized retouching).
\end{itemize}

The separation of generated and edited classes enables the classifier to learn
distinct signatures: global spectral fingerprints often indicate fully
generated imagery, while localized residual anomalies better characterise
editing operations.

\subsection{Dataset Size, Splits and Export}
Recorded counts are: real = 26,000; ai_generated = 26,000; ai_edited = 25,865; total = 77,865. The pipeline writes exported splits to export directories. While nominal YAML-configured ratios exist (70/15/15), the implemented cluster-aware splitter produces an empirical export approximating 40/30/30 in some runs; the split implementation is in the splitter module and is orchestrated by the dataset pipeline.

\subsection{Class Balancing and Quota Sampling}
The production pipeline enforces per-source quotas and overall class balance to
prevent dominance of any single class or source. The sampling logic is in
The sampling logic is implemented in the sampler module and deterministic
behaviour is ensured via \texttt{set_global_seed} in the dataset pipeline.
Balanced datasets reduce training bias and ensure that per-class metrics are
meaningful for operational evaluation.

\subsection{Source Diversity and Manipulation Types}
The dataset aggregates diverse capture conditions and manipulation methods to
improve generalisation. Diversity includes multiple camera models, compression
artifacts, and generative model variants; manipulation types include full
generation, inpainting, and face-swapping. The pipeline preserves provenance
metadata (for example \path{export_index.csv}) to support stratified analyses
by source or manipulation family.

\subsection{Deduplication, Validation and Dataset Quality}
Deduplication is performed with perceptual hashing and validation stages check
resolution, blur and metadata (see the deduplicator and validator modules).
Cluster-aware splitting
ensures that related items (e.g. frames extracted from the same video) remain
in the same split to prevent leakage. These steps are critical because low-
quality or duplicated training examples encourage memorisation of source
artifacts rather than learning transferable forensic cues.

\section{Data Preprocessing}
\subsection{Overview}
Preprocessing responsibilities are split across dataset-level filtering and
model-input preparation. Model-input preprocessing is implemented primarily in
the preprocessing module and uses Albumentations for training-time
augmentations and Torchvision transforms for deterministic inference-time
processing. Optional preprocessing modules (\texttt{SRMLayer}, \texttt{FFTLayer})
are implemented in dedicated SRM and FFT modules and are applied via the
\texttt{PreprocessNet} wrapper in the training module.

\subsection{Image Loading and Channel Ordering}
At data ingest, images are read with OpenCV (BGR format) within the dataloader
module. The loader performs an explicit BGR→RGB conversion prior to
augmentation to align with pretrained backbones' expected channel ordering and
avoid inadvertent color-space anomalies.

\subsection{Resizing and Normalization}
All inputs are resized to 224×224 during augmentation. Pixel normalisation uses
ImageNet mean and standard deviation (mean=[0.485,0.456,0.406],
std=[0.229,0.224,0.225]) to maintain consistency with pretrained feature
initialisations. Conversion to PyTorch channel ordering (C,H,W) and floating
tensor is performed by \texttt{ToTensorV2} as the final augmentation step.

\subsection{Augmentation Modes and Techniques}
Training augmentations are exposed as three named transforms via
\texttt{get_train_transform(mode)}:
\begin{itemize}
  \item \texttt{light}: Resize(224,224), HorizontalFlip(p=0.5), Normalize(),
  \texttt{ToTensorV2}.
  \item \texttt{standard} (default): Resize(224,224), HorizontalFlip(p=0.5),
  Rotate(limit=15,p=0.5), RandomBrightnessContrast(p=0.5),
  ImageCompression(quality_range=(30,100),p=0.4), Normalize(), \texttt{ToTensorV2}.
  \item \texttt{strong}: as \texttt{standard} with intensified ImageCompression
  (20–95), additive noise (GaussNoise or ISONoise) and light blur (Gaussian or
  Median) to simulate harsh capture conditions.
\end{itemize}

Rationale: Compression simulation addresses lossy storage/transmission; noise
injection models sensor or post-processing noise; geometric transforms improve
robustness to framing variations; blur reduces sensitivity to fine-grained
texture so models learn coarser but more generalisable cues.

\subsection{Training vs Inference Preprocessing}
Training uses stochastic Albumentations to promote invariance, while
inference applies deterministic Torchvision transforms (Resize, CenterCrop,
ToTensor, Normalize) to ensure repeatable inputs and stable evaluation.

\section{Feature Extraction}
\subsection{Candidate Backbones and Feature Dimensions}
The codebase supports multiple backbone families (constructed in
\path{scripts/training/train_full.py}): ResNet-18/ResNet-50 (final features
512/2048), ConvNeXt-Tiny/Small (final features 768), EfficientNet-B3 and
ViT-B/16 (supported with special handling for token pooling). Using
ImageNet-pretrained weights provides robust low-level filters and accelerates
convergence when fine-tuning for forensic tasks.

\subsection{Convolutional Inductive Biases for Forensics}
Convolutional architectures are well-suited for forensic analysis due to their
local receptive fields and hierarchical composition. Early layers capture
edges and textures while deeper layers encode increasingly semantic
representations; this hierarchy supports detection of both local manipulations
(inpainting seams) and more global inconsistencies (colour mismatches or
artificial textures). Transformer-based backbones are supported to compare
their global attention properties, but their tokenisation implies different
trade-offs.

\subsection{Pooling and Classification Head}
After spatial aggregation (GeM or average pooling) the pooled vector is
subjected to dropout (default 0.4) and then passed to a linear classification
head mapping to three logits. The head construction is performed by
\texttt{make_head(in_features)} in \path{scripts/training/train_full.py} and
is tuned to each backbone's final feature dimensionality.

\section{Preprocessing Feature Modules}
\subsection{\textbf{Steganalysis Rich Model (SRM)}}
\textbf{Implementation:} \texttt{SRMLayer} (\path{scripts/preprocessing/srm.py})
computes three fixed high-pass residual channels using handcrafted kernels
(Laplacian 5×5, horizontal 5×5, vertical 5×5) applied to a grayscale image;
residuals are clamped (default clamp range 2.0) and concatenated with RGB to
form a 6-channel input. The helper \texttt{adapt_conv1_for_srm} modifies the
backbone's first convolution to accept 6 channels by initialising additional
weights to 0.1× the pretrained RGB weights.

Justification: SRM residuals amplify small-scale discontinuities and blending
seams that are often obscured by colour and texture in the RGB domain. By
presenting these residuals explicitly to the network, early convolutions can
detect structural anomalies that are indicative of manipulation.

\subsection{\textbf{Fast Fourier Transform (FFT)}}
\textbf{Implementation:} \texttt{FFTLayer} (\path{scripts/preprocessing/fft.py})
computes the log-magnitude of the centred 2D Fourier transform of the
grayscale image, normalises it to [0,1], and returns a B×1×H×W tensor. When
enabled, this channel is concatenated after RGB and SRM channels (rgb →
residuals → fft) and the backbone's first convolution is adapted accordingly.

Justification: Frequency-domain analysis reveals global periodicities and
power-spectrum anomalies often characteristic of generative models or
resampling introduced by editing tools. FFT magnitude provides complementary
evidence to SRM and RGB features and can distinguish manipulation classes
that are ambiguous in the spatial domain.

\section{Attention and Pooling Mechanisms}
\subsection{\textbf{CBAM Attention}: Channel and Spatial Attention}
CBAM is implemented in \path{scripts/modules/attention_heads.py} and applies
channel attention (via pooled statistics and an MLP) followed by spatial
attention (via concatenated avg/max maps and a convolution). Applied to late
feature maps, CBAM reweights channels and spatial locations to emphasise
forensically salient information.

\subsection{\textbf{GeM Pooling}}
\textbf{Implementation:} \texttt{GeM} (\path{scripts/modules/attention_heads.py}).
Given activation map $X \in \mathbb{R}^{C\times H\times W}$ GeM computes:
\[
\mathrm{GeM}(X)_c =
\left(
\frac{1}{HW}
\sum_{i=1}^{H}
\sum_{j=1}^{W}
X_{c,i,j}^{p}
\right)^{1/p}
\]
The parameter $p$ (default 3.0) is registered as a tensor and may be frozen
or learned. GeM interpolates between average ($p=1$) and max pooling
($p\to\infty$), giving stronger emphasis to local high activations which are
useful when small regions contain manipulation evidence.

\subsection{Role of Attention and GeM in Forensics}
Attention reduces the influence of irrelevant background structure and lets the
model focus on high-signal regions. GeM ensures that strong local cues are
preserved during aggregation rather than being diluted by large spatial
averages; this property is particularly beneficial when manipulations occupy a
small fraction of the image.

\section{Model Architecture}
\subsection{Input and \texttt{PreprocessNet} Wrapper}
The input pipeline optionally applies \texttt{SRMLayer} and/or \texttt{FFTLayer}
to an RGB tensor and concatenates channels in the order rgb → residuals → fft.
The concatenated tensor is provided to an adapted backbone via the
\texttt{PreprocessNet} wrapper implemented in \path{scripts/training/train_full.py}.

\subsection{Backbone, Attention and Pooling}
The adapted backbone (ResNet/ConvNeXt/EfficientNet/ViT) computes hierarchical
features. \textbf{CBAM Attention} blocks may be inserted after late stages
(for ResNet, after \path{layer4}; for ConvNeXt, after \path{features}) and
\textbf{GeM Pooling} may replace the global average pooling layer. The
resulting pooled vector is passed through a dropout layer and a linear
classification head (softmax for probabilities).

\subsection{Step-by-step Data Flow}
\begin{enumerate}[label=\textbf{Step \arabic*:}]
  \item Image file loaded (cv2) and converted BGR→RGB.
  \item Training augmentations applied (Resize → geometric → photometric →
  Normalize → \texttt{ToTensorV2}) or deterministic inference transforms used.
  \item \texttt{SRMLayer} / \texttt{FFTLayer} modules compute additional
  channels if enabled and outputs are concatenated with RGB.
  \item The concatenated tensor is forwarded through the adapted backbone, with
  optional \textbf{CBAM Attention} applied to late feature maps.
  \item \textbf{GeM Pooling} or average pooling aggregates spatial maps to a
  vector, followed by dropout and the final linear head producing three logits.
\end{enumerate}

\section{Training Procedure}
\subsection{Losses and Class-weighting}
The training code implements CrossEntropy, class-weighted CrossEntropy,
\textbf{Focal Loss} (\texttt{FocalLoss}) and a factory \texttt{build_criterion}
for WeightedFocalLoss. Default class weights are \path{[1.5, 1.0, 1.5]}. The
default focal gamma in code is 3.0; however model cards and many experiments
use $\gamma=2.0$ in practice. Label smoothing (commonly 0.1) is supported to
reduce overconfidence.

Rationale: Focal loss focuses optimisation on hard examples, which is useful
when manipulated examples are subtle. Class weights compensate for residual
sampling imbalance and emphasize operationally critical classes.

\subsection{Optimisers, Schedules and Hyperparameters}
Optimisers: Adam (default) and SGD (option). Default weight decay: 1e-4. The
learning rate default is 1e-4. Schedulers: \textbf{Cosine Annealing}
restarts (CosineAnnealingWarmRestarts with \texttt{T_0=20}, \texttt{T_mult=2}) for \texttt{lr_schedule='cosine'}, or ReduceLROnPlateau for \texttt{lr_schedule='plateau'}.
Batch size default is 64 and default epochs is 50 (baseline scripts use
smaller defaults for quick iterations). Early stopping and checkpointing save
the best model (for example to \path{best_model.pth}) and per-epoch metrics
are written to \path{metrics.csv} and \path{training_summary.json}.

Rationale: Cosine restarts provide periodic learning rate resets that help
escape local minima during long runs; Adam provides adaptive gradient steps that
are robust to noisy gradients introduced by augmentations.

\subsection{Mixed Precision and Performance}
Training uses AMP (\texttt{torch.amp.GradScaler}, \texttt{autocast}) and
enables \texttt{torch.compile} when available. These pragmatics increase
throughput and reduce memory footprint, enabling larger batch sizes and
faster experimentation.

\section{Evaluation Metrics}
\subsection{Metrics Implemented}
\begin{itemize}
  \item \textbf{Accuracy}: overall fraction correct; logged as
  \path{train_acc} and \path{val_acc}.
  \item \textbf{Precision / Recall}: computed per-class; precision measures the
  correctness of positive predictions while recall measures completeness of
  detection for the true positives.
  \item \textbf{F1 score}: harmonic mean of precision and recall; code computes
  macro and per-class F1 using \path{sklearn.metrics.f1_score}.
  \item \textbf{Confusion matrix}: reconstructed from saved logits and labels to
  provide detailed misclassification patterns.
\end{itemize}

\subsection{ROC-AUC and Offline Analysis}
ROC-AUC is not computed by default within the training scripts but can be
derived offline from saved logits and labels. ROC-AUC is threshold-independent
and especially useful when selecting operational thresholds that balance false
positives and false-negatives.

\subsection{Complementary Metric Usage}
Accuracy gives a single-number summary but can conceal class-specific failures;
precision, recall and F1 provide per-class diagnostic insight, while the
confusion matrix reveals asymmetric confusions that guide dataset curation and
loss reweighting.

\section{Methodology Summary}
This chapter has presented a principled methodology integrating dataset
engineering, robust preprocessing, multi-domain feature extraction, attention
and pooling mechanisms, and carefully chosen optimisation strategies. The
approach is implemented in the repository with explicit modules to ensure
reproducibility. Each design decision is motivated by forensic considerations:
\textbf{Steganalysis Rich Model (SRM)} emphasises high-frequency residuals for
local edit detection; \textbf{Fast Fourier Transform (FFT)} reveals global
spectral fingerprints; \textbf{GeM Pooling} and \textbf{CBAM Attention}
concentrate model capacity on salient, high-signal regions; and augmentation
plus class-aware sampling promote robust generalisation. The subsequent
chapters evaluate this methodology empirically and discuss limitations and
avenues for future work.

\section{Diagrams}
The following subsections describe recommended diagrams. Each subsection lists
\textbf{Components} and a short \textbf{Purpose} note.

\subsection{Overall System Pipeline}
\textbf{Components}
\begin{itemize}
  \item raw sources
  \item indexing
  \item validation
  \item deduplication
  \item sampling / balancing
  \item cluster-aware splitter
  \item exported train / val / test
  \item training (\texttt{PreprocessNet} → Backbone → Head)
  \item evaluation
  \item model cards / metrics
\end{itemize}
\textbf{Purpose}
End-to-end reproducible workflow from raw data to published experiment artifacts.
 
\subsection{Dataset Preprocessing Pipeline}
\textbf{Components}
\begin{itemize}
  \item indexing
  \item validation checks (resolution / blur)
  \item deduplication (pHash)
  \item quota sampling
  \item cluster-aware split
  \item export
\end{itemize}
\textbf{Purpose}
Clarifies steps that ensure dataset quality and avoid leakage; show conditional branches (invalid → rejected) and deterministic seeding.

\subsection{CNN Backbone Architecture}
\textbf{Components}
\begin{itemize}
  \item input
  \item stem convs
  \item residual / ConvNeXt blocks
  \item stage outputs
  \item pooling (Avg / GeM)
  \item dropout
  \item FC head
\end{itemize}
\textbf{Purpose}
Hierarchical feature extraction for forensic signals (feature sizes: ResNet-18: 512; ConvNeXt: 768; ResNet-50: 2048).

\subsection{SRM Filter Integration}
\textbf{Components}
\begin{itemize}
  \item RGB → grayscale
  \item three SRM kernels (Laplacian, horizontal, vertical)
  \item residual maps
  \item concatenation with RGB
  \item adapted conv1 (0.1× init for extra channels)
\end{itemize}
\textbf{Purpose}
Shows how residual channels are fused at model input and highlights how SRM buffers are stored in checkpoints.

\subsection{FFT Feature Extraction}
\textbf{Components}
\begin{itemize}
  \item grayscale → FFT → fftshift → log-magnitude
  \item normalise → FFT channel concatenated with RGB / SRM
\end{itemize}
\textbf{Purpose}
Documents spectral fingerprint extraction and the transformation from spatial to frequency domain and back to a spatially-aligned magnitude image.

\subsection{Attention Module (CBAM)}
\textbf{Components}
\begin{itemize}
  \item input feature map
  \item channel attention (avg/max pools → MLP → sigmoid)
  \item multiply
  \item spatial attention (concat avg/max → conv → sigmoid)
  \item multiply
\end{itemize}
\textbf{Purpose}
Describes sequential application (channel → spatial) and how attention reweights channels and locations.

\subsection{Feature Extraction Pipeline}
\textbf{Components}
\begin{itemize}
  \item RGB → (optional SRM / FFT)
  \item adapted backbone
  \item attention
  \item pooling (GeM)
  \item feature vector → classifier head
\end{itemize}
\textbf{Purpose}
Shows multi-domain fusion and final prediction flow.

\subsection{Full Model Architecture}
\textbf{Components}
\begin{itemize}
  \item \texttt{PreprocessNet} wrapper
  \item backbone
  \item optional \textbf{CBAM Attention}
  \item \textbf{GeM Pooling}
  \item dropout
  \item fully connected classifier
\end{itemize}
\textbf{Purpose}
Full neural network inference pipeline with tensor dimension and feature flow annotations.

\subsection{Training Workflow}
\textbf{Components}
\begin{itemize}
  \item DataLoader
  \item augmentations
  \item forward pass
  \item loss computation
  \item backward pass
  \item optimizer step
  \item scheduler update
  \item checkpoint saving
\end{itemize}
\textbf{Purpose}
Describes the training loop execution.

\subsection{Evaluation Pipeline}
\textbf{Components}
\begin{itemize}
  \item test dataset
  \item inference
  \item logits collection
  \item metrics computation (accuracy, precision, recall, F1)
  \item confusion matrix generation
\end{itemize}
\textbf{Purpose}
Evaluation and result reporting.

\end{document}
